% Part: first-order-logic
% Chapter: sequent-calculus
% Section: soundness.tex

\documentclass[../../../include/open-logic-section]{subfiles}

\begin{document}

\iftag{FOL}
      {\olfileid{fol}{seq}{sou}}
      {\olfileid{pl}{seq}{sou}}

\olsection{Correctheid}

\begin{explain}
!!^a{derivation}ssysteem, zoals de sequentenrekening, is \emph{correct}
als het geen zaken kan !!{derive} die in werkelijkheid niet gelden.
Correctheid is dus een soort gegarandeerde veiligheidseigenschap van
!!{derivation}ssystemen. Afhankelijk van de bewijstheoretische
eigenschap waar het om gaat, willen we bijvoorbeeld weten dat
\begin{enumerate}
\item elke !!{derivable}~$!A$ \iftag{FOL}{geldig}{een tautologie} is;
\item als !!a{sentence} uit enkele andere zinnen !!{derivable}
  is, die zin ook een gevolg daarvan is;
\item als een verzameling !!{sentence}s inconsistent is, ze
  onvervulbaar is.
\end{enumerate}
Dit zijn belangrijke eigenschappen van !!a{derivation}ssysteem. Als
een ervan niet geldt, schiet het !!{derivation}ssysteem tekort: het zou
te veel !!{derive}. Het is daarom van het grootste belang om de
correctheid van een !!{derivation}ssysteem vast te stellen.

Omdat al deze bewijstheoretische eigenschappen via de
!!{derivability} van bepaalde sequenten in de sequentenrekening zijn
gedefinieerd, vereist een bewijs van (1)--(3) dat we iets bewijzen over
de semantische eigenschappen van !!{derivable} sequenten. We zullen
eerst definiëren wat het betekent dat een sequent \emph{geldig} is en
vervolgens aantonen dat elke !!{derivable} sequent geldig is. (1)--(3)
volgen dan als corollaria uit dit resultaat.
\end{explain}

\begin{defn}
\iftag{FOL}{!!^a{structure}~$\Struct
  M$}{!!^a{valuation}~$\pAssign{v}$} \emph{vervult} een sequent
$\Gamma \Sequent \Delta$ dan en slechts dan als ofwel
$\iftag{FOL}{\Sat/{M}}{\pSat/{v}}{!A}$ voor een $!A \in \Gamma$, ofwel
$\iftag{FOL}{\Sat{M}}{\pSat{v}}{!A}$ voor een $!A \in \Delta$.

Een sequent is \emph{geldig} dan en slechts dan als elke
\iftag{FOL}{!!{structure}~$\Struct M$}{!!{valuation}~$\pAssign{v}$}
hem vervult.
\end{defn}

\begin{thm}[Correctheid]
\ollabel{thm:sequent-soundness} Als $\Log{LK}$ de sequent $\Theta
\Sequent \Xi$ !!{derive}, dan is $\Theta \Sequent \Xi$ geldig.
\end{thm}

\begin{proof}
Laat $\pi$ !!a{derivation} van $\Theta \Sequent \Xi$ zijn. We passen
inductie toe naar het aantal afleidingsstappen~$n$ in~$\pi$.

Als het aantal afleidingsstappen~$0$ is, bestaat $\pi$ alleen uit een
beginsequent. Elke beginsequent $!A \Sequent !A$ is klaarblijkelijk
geldig, want voor elke \iftag{FOL}{$\Struct M$}{$\pAssign{v}$} geldt ofwel
$\iftag{FOL}{\Sat/{M}}{\pSat/{v}}{!A}$ of
$\iftag{FOL}{\Sat{M}}{\pSat{v}}{!A}$.

Als het aantal afleidingsstappen groter is dan~0, onderscheiden we
gevallen naar het type van de onderste afleidingsstap. Volgens de
inductiehypothese mogen we aannemen dat de premissen van die stap
geldig zijn, want de !!{derivation} van elke premisse bevat minder dan
$n$ afleidingsstappen.

Eerst behandelen we de mogelijke afleidingsstappen met slechts één
premisse.
\begin{enumerate}
\item De laatste afleidingsstap is een verzwakking. Dan is $\Theta
  \Sequent \Xi$ ofwel $!A, \Gamma \Sequent \Delta$ (als de laatste
  stap \LeftR{\Weakening} is), ofwel $\Gamma \Sequent \Delta, !A$
  (als die stap \RightR{\Weakening} is), en eindigt de !!{derivation}
  in een van
  \begin{prooftree}
    \AxiomC{}
    \Deduce$\Gamma \fCenter \Delta$
    \RightLabel{\LeftR{\Weakening}}
    \UnaryInf$!A, \Gamma \fCenter \Delta$
    \DisplayProof\qquad\bottomAlignProof
    \AxiomC{}
    \Deduce$\Gamma \fCenter \Delta$
    \RightLabel{\RightR{\Weakening}}
    \UnaryInf$\Gamma \fCenter \Delta, !A$
  \end{prooftree}
  Volgens de inductiehypothese is $\Gamma \Sequent \Delta$ geldig, dat
  wil zeggen: voor elke
  \iftag{FOL}{!!{structure}~$\Struct{M}$}{!!{valuation}~$\pAssign{v}$},
  is er ofwel een $!C \in \Gamma$ waarvoor
  $\iftag{FOL}{\Sat/{M}}{\pSat/{v}}{!C}$, ofwel een $!C \in \Delta$
  waarvoor $\iftag{FOL}{\Sat{M}}{\pSat{v}}{!C}$.
  
  Als $\iftag{FOL}{\Sat/{M}}{\pSat/{v}}{!C}$ voor een $!C \in
  \Gamma$, dan geldt $!C \in \Theta$, want $\Theta = !A,
  \Gamma$. Er is dus een $!C \in \Theta$ waarvoor
  $\iftag{FOL}{\Sat/{M}}{\pSat/{v}}{!C}$. Evenzo geldt: als
  $\iftag{FOL}{\Sat{M}}{\pSat{v}}{!C}$ voor een $!C \in \Delta$, dan
  geldt $!C \in \Xi$ en bovendien
  $\iftag{FOL}{\Sat{M}}{\pSat{v}}{!C}$ voor een $!C \in \Xi$.
  Bijgevolg is $\Theta \Sequent \Xi$ geldig.
\item De laatste afleidingsstap is \LeftR{\lnot}. De premisse van die
  stap is dan $\Gamma \Sequent \Delta, !A$ en de conclusie is $\lnot
  !A, \Gamma \Sequent \Delta$. De !!{derivation} eindigt dus in
  \begin{prooftree}
    \AxiomC{}
    \Deduce$\Gamma \fCenter \Delta, !A$
    \RightLabel{\LeftR{\lnot}}
    \UnaryInf$\lnot !A, \Gamma \fCenter \Delta$
  \end{prooftree}
  waarbij $\Theta = \lnot !A, \Gamma$ en $\Xi = \Delta$.

  Volgens de inductiehypothese is $\Gamma \Sequent \Delta, !A$
  geldig. Voor elke \iftag{FOL}{$\Struct{M}$}{$\pAssign{v}$} geldt dus
  ofwel (a) $\iftag{FOL}{\Sat/{M}}{\pSat/{v}}{!C}$ voor een $!C \in
  \Gamma$, ofwel (b) $\iftag{FOL}{\Sat{M}}{\pSat{v}}{!C}$ voor een
  $!C \in \Delta$, ofwel (c)
  $\iftag{FOL}{\Sat{M}}{\pSat{v}}{!A}$. We willen aantonen dat ook
  $\Theta \Sequent \Xi$ geldig is. Laat
  \iftag{FOL}{$\Struct{M}$ !!a{structure} zijn}{$\pAssign{v}$
    !!a{valuation} zijn}. Als (a) geldt, is er een $!C \in \Gamma$
  waarvoor $\iftag{FOL}{\Sat/{M}}{\pSat/{v}}{!C}$, en geldt
  $!C \in \Theta$. Als (b) geldt, is er een $!C \in \Delta$ waarvoor
  $\iftag{FOL}{\Sat{M}}{\pSat{v}}{!C}$, en geldt $!C \in \Xi$. Als ten
  slotte $\iftag{FOL}{\Sat{M}}{\pSat{v}}{!A}$, dan
  $\iftag{FOL}{\Sat/{M}}{\pSat/{v}}{\lnot !A}$. Omdat $\lnot !A \in
  \Theta$, is er een $!C \in \Theta$ waarvoor
  $\iftag{FOL}{\Sat/{M}}{\pSat/{v}}{!C}$. Bijgevolg is $\Theta
  \Sequent \Xi$ geldig.
\item De laatste afleidingsstap is \RightR{\lnot}: opgave.
\item De laatste afleidingsstap is \LeftR{\land}. Er zijn twee
  varianten: $!A \land !B$ kan links worden afgeleid uit $!A$ of uit
  $!B$ aan de linkerkant van de premisse. In het eerste geval eindigt
  $\pi$ in
  \begin{prooftree}
    \AxiomC{}
    \Deduce$!A, \Gamma \fCenter \Delta$
    \RightLabel{\LeftR{\land}}
    \UnaryInf$!A \land !B, \Gamma \fCenter \Delta$
  \end{prooftree}
  en is $\Theta = !A \land !B, \Gamma$, terwijl $\Xi = \Delta$.
  Beschouw een
  \iftag{FOL}{!!a{structure}~$\Struct
    M$}{!!a{valuation}~$\pAssign{v}$}. Volgens de inductiehypothese is
  $!A, \Gamma \Sequent \Delta$ geldig, zodat ofwel (a)
  $\iftag{FOL}{\Sat/{M}}{\pSat/{v}}{!A}$, (b)
  $\iftag{FOL}{\Sat/{M}}{\pSat/{v}}{!C}$ voor een $!C \in \Gamma$,
  of (c)~$\iftag{FOL}{\Sat{M}}{\pSat{v}}{!C}$ voor een $!C \in
  \Delta$. In geval (a) geldt
  $\iftag{FOL}{\Sat/{M}}{\pSat/{v}}{!A \land !B}$. Er is dus een $!C
  \in \Theta$ (namelijk $!A \land !B$) waarvoor
  $\iftag{FOL}{\Sat/{M}}{\pSat/{v}}{!C}$. In geval (b) is er een $!C
  \in \Gamma$ waarvoor $\iftag{FOL}{\Sat/{M}}{\pSat/{v}}{!C}$, en
  geldt $!C \in \Theta$. In geval (c) is er een $!C \in
  \Delta$ waarvoor $\iftag{FOL}{\Sat{M}}{\pSat{v}}{!C}$, en geldt
  $!C \in \Xi$, want $\Xi = \Delta$. In elk geval vervult
  \iftag{FOL}{$\Struct M$}{$\pAssign{v}$} dus $!A \land !B, \Gamma
  \Sequent \Delta$. Omdat \iftag{FOL}{$\Struct{M}$}{$\pAssign{v}$}
  willekeurig was, is $\Gamma \Sequent \Delta$ geldig. Het geval
  waarin $!A \land !B$ uit $!B$ wordt afgeleid, verloopt hetzelfde,
  met $!B$ in plaats van $!A$.
\item De laatste afleidingsstap is \RightR{\lor}. Er zijn twee
  varianten: $!A \lor !B$ kan rechts worden afgeleid uit $!A$ of uit
  $!B$ aan de rechterkant van de premisse. In het eerste geval eindigt
  $\pi$ in
  \begin{prooftree}
    \AxiomC{}
    \Deduce$\Gamma \fCenter \Delta, !A$
    \RightLabel{\RightR{\lor}}
    \UnaryInf$\Gamma \fCenter \Delta, !A \lor !B$
  \end{prooftree}
  Nu is $\Theta = \Gamma$ en $\Xi = \Delta, !A \lor !B$. Beschouw een
  \iftag{FOL}{!!a{structure}~$\Struct{M}$}{!!a{valuation}~$\pAssign{v}$}.
  Omdat $\Gamma \Sequent \Delta, !A$ geldig is, geldt ofwel (a)
  $\iftag{FOL}{\Sat{M}}{\pSat{v}}{!A}$, (b)
  $\iftag{FOL}{\Sat/{M}}{\pSat/{v}}{!C}$ voor een $!C \in \Gamma$,
  of (c)~$\iftag{FOL}{\Sat{M}}{\pSat{v}}{!C}$ voor een $!C \in
  \Delta$. In geval (a) geldt
  $\iftag{FOL}{\Sat{M}}{\pSat{v}}{!A \lor !B}$. In geval (b) is er een
  $!C \in \Gamma$ waarvoor $\iftag{FOL}{\Sat/{M}}{\pSat/{v}}{!C}$.
  In geval (c) is er een $!C \in \Delta$ waarvoor
  $\iftag{FOL}{\Sat{M}}{\pSat{v}}{!C}$. In elk geval vervult
  \iftag{FOL}{$\Struct{M}$}{$\pAssign{v}$} dus $\Gamma \Sequent
  \Delta, !A \lor !B$, oftewel $\Theta \Sequent \Xi$. Omdat
  \iftag{FOL}{$\Struct{M}$}{$\pAssign{v}$} willekeurig was, is $\Theta
  \Sequent \Xi$ geldig. Het geval waarin $!A \lor !B$ uit $!B$ wordt
  afgeleid, verloopt hetzelfde, met $!B$ in plaats van $!A$.
\item De laatste afleidingsstap is \RightR{\lif}. Dan eindigt $\pi$ in
  \begin{prooftree}
    \AxiomC{}
    \Deduce$!A, \Gamma \fCenter \Delta, !B$
    \RightLabel{\RightR{\lif}}
    \UnaryInf$\Gamma \fCenter \Delta, !A \lif !B$
  \end{prooftree}
  Ook nu zegt de inductiehypothese dat de premisse geldig is; we willen
  aantonen dat de conclusie eveneens geldig is. Laat
  \iftag{FOL}{$\Struct{M}$}{$\pAssign{v}$} willekeurig zijn. Omdat
  $!A, \Gamma \Sequent \Delta, !B$ geldig is, doet zich ten minste een
  van de volgende gevallen voor: (a)
  $\iftag{FOL}{\Sat/{M}}{\pSat/{v}}{!A}$, (b)
  $\iftag{FOL}{\Sat{M}}{\pSat{v}}{!B}$, (c)
  $\iftag{FOL}{\Sat/{M}}{\pSat/{v}}{!C}$ voor een~$~!C \in \Gamma$,
  of (d) $\iftag{FOL}{\Sat{M}}{\pSat{v}}{!C}$ voor een $!C \in
  \Delta$. In de gevallen (a) en (b) geldt
  $\iftag{FOL}{\Sat{M}}{\pSat{v}}{!A \lif !B}$, zodat er een $!C \in
  \Delta, !A \lif !B$ is waarvoor
  $\iftag{FOL}{\Sat{M}}{\pSat{v}}{!C}$. In geval (c) geldt voor een
  $!C \in \Gamma$ dat $\iftag{FOL}{\Sat/{M}}{\pSat/{v}}{!C}$. In
  geval (d) geldt voor een $!C \in \Delta$ dat
  $\iftag{FOL}{\Sat{M}}{\pSat{v}}{!C}$. In elk geval vervult
  \iftag{FOL}{$\Struct{M}$}{$\pAssign{v}$} de sequent $\Gamma \Sequent
  \Delta, !A \lif !B$. Omdat \iftag{FOL}{$\Struct{M}$}{$\pAssign{v}$}
  willekeurig was, is $\Gamma \Sequent \Delta, !A \lif !B$ geldig.
  \iftag{FOL}{%
\item De laatste afleidingsstap is \LeftR{\lforall}. Dan zijn er
  !!a{formula}~$!A(x)$ en een gesloten term~$t$ zodanig dat $\pi$
  eindigt in
  \begin{prooftree}
    \AxiomC{}
    \Deduce$!A(t), \Gamma \fCenter \Delta$
    \RightLabel{\LeftR{\lforall}}
    \UnaryInf$\lforall[x][!A(x)], \Gamma \fCenter \Delta$
  \end{prooftree}
  We willen aantonen dat de conclusie $\lforall[x][!A(x)], \Gamma
  \Sequent \Delta$ geldig is. Beschouw !!a{structure}~$\Struct M$.
  Omdat de premisse $!A(t), \Gamma \Sequent \Delta$ geldig is, geldt
  ofwel (a) $\Sat/{M}{!A(t)}$, ofwel (b) $\Sat/{M}{!C}$ voor een $!C
  \in \Gamma$, ofwel (c)~$\Sat{M}{!C}$ voor een $!C \in \Delta$. In
  geval (a) geldt volgens \olref[syn][sem]{prop:quant-terms}: als
  $\Sat{M}{\lforall[x][!A(x)]}$, dan $\Sat{M}{!A(t)}$. Omdat
  $\Sat/{M}{!A(t)}$, geldt $\Sat/{M}{\lforall[x][!A(x)]}$. In de
  gevallen (b) en (c) vervult $\Struct{M}$ eveneens $\lforall[x][!A(x)],
  \Gamma \Sequent \Delta$. Omdat $\Struct M$ willekeurig was, is
  $\lforall[x][!A(x)], \Gamma \Sequent \Delta$ geldig.
\item De laatste afleidingsstap is \RightR{\lexists}: opgave.
\item De laatste afleidingsstap is \RightR{\lforall}. Dan zijn er
  !!a{formula}~$!A(x)$ en !!a{constant}~$a$ zodanig dat $\pi$ eindigt in
  \begin{prooftree}
    \AxiomC{}
    \Deduce$\Gamma \fCenter \Delta, !A(a)$
    \RightLabel{\RightR{\lforall}}
    \UnaryInf$\Gamma \fCenter \Delta, \lforall[x][!A(x)]$
  \end{prooftree}
  waarbij aan de eigenvariabelevoorwaarde is voldaan: $a$ komt niet
  voor in $!A(x)$, $\Gamma$ of $\Delta$. Volgens de inductiehypothese
  is de premisse van de laatste afleidingsstap geldig. We moeten
  aantonen dat ook de conclusie geldig is, dus dat voor elke
  !!{structure}~$\Struct M$ ofwel (a)
  $\Sat{M}{\lforall[x][!A(x)]}$, ofwel (b) $\Sat/{M}{!C}$ voor een $!C
  \in \Gamma$, ofwel (c)~$\Sat{M}{!C}$ voor een $!C \in \Delta$.

  Zij $\Struct{M}$ een willekeurige !!{structure}. Als (b) of (c)
  geldt, zijn we klaar. Neem dus aan dat geen van beide geldt: voor
  alle $!C \in \Gamma$ geldt $\Sat{M}{!C}$, en voor alle $!C \in
  \Delta$ geldt $\Sat/{M}{!C}$. We moeten aantonen dat (a) geldt, dus
  $\Sat{M}{\lforall[x][!A(x)]}$. Volgens
  \olref[syn][ass]{prop:sat-quant} volstaat het aan te tonen dat
  $\Sat{M}{!A(x)}[s]$ voor alle bedelingen~$s$. Laat $s$ dus een
  willekeurige bedeling zijn. Beschouw de structuur~$\Struct{M'}$ die
  alleen daarin van~$\Struct{M}$ verschilt dat $\Assign{a}{M'} =
  s(x)$. Volgens \olref[syn][ext]{cor:extensionality-sent} geldt voor
  elke $!C \in \Gamma$ dat $\Sat{M'}{!C}$, omdat $a$ niet in~$\Gamma$
  voorkomt, en voor elke $!C \in \Delta$ dat $\Sat/{M'}{!C}$. De
  premisse is echter geldig, dus $\Sat{M'}{!A(a)}$. Volgens
  \olref[syn][ass]{prop:sentence-sat-true} geldt
  $\Sat{M'}{!A(a)}[s]$, want $!A(a)$ is een zin. Nu
  $\varAssign{s}{s}{x}$ met $s(x) = \Value{a}{M'}[s]$, omdat we
  $\Struct{M'}$ precies zo hebben gedefinieerd. We kunnen dus
  \olref[syn][ext]{prop:ext-formulas} toepassen en verkrijgen
  $\Sat{M'}{!A(x)}[s]$. Omdat $a$ niet in~$!A(x)$ voorkomt, volgt uit
  \olref[syn][ext]{prop:extensionality} dat
  $\Sat{M}{!A(x)}[s]$. Aangezien $s$ willekeurig was, hebben we bewezen
  dat $\Sat{M}{!A(x)}[s]$ voor alle bedelingen.
\item De laatste afleidingsstap is \LeftR{\lexists}: opgave.
}{}
\end{enumerate}
Nu behandelen we de mogelijke afleidingsstappen met twee premissen.
\begin{enumerate}
\item De laatste afleidingsstap is een snede. Dan eindigt $\pi$ in
  \begin{prooftree}
    \AxiomC{}
    \Deduce$\Gamma \fCenter \Delta, !A$
    \AxiomC{}
    \Deduce$!A, \Pi \fCenter \Lambda$
    \RightLabel{\Cut}
    \BinaryInf$\Gamma, \Pi \fCenter \Delta, \Lambda$
  \end{prooftree}
  Laat \iftag{FOL}{$\Struct{M}$ !!a{structure} zijn}{$\pAssign{v}$
    !!a{valuation} zijn}. Volgens de inductiehypothese zijn de
  premissen geldig, zodat \iftag{FOL}{$\Struct{M}$}{$\pAssign{v}$}
  beide premissen vervult. We onderscheiden twee gevallen: (a)
  $\iftag{FOL}{\Sat/{M}}{\pSat/{v}}{!A}$ en (b)
  $\iftag{FOL}{\Sat{M}}{\pSat{v}}{!A}$. In geval (a) kan
  \iftag{FOL}{$\Struct{M}$}{$\pAssign{v}$} de linkerpremisse alleen
  vervullen als hij $\Gamma \Sequent \Delta$ vervult. Dan vervult hij
  ook de conclusie. In geval (b) kan
  \iftag{FOL}{$\Struct{M}$}{$\pAssign{v}$} de rechterpremisse alleen
  vervullen als hij $\Pi \setminus \Lambda$ vervult. Ook dan vervult
  \iftag{FOL}{$\Struct{M}$}{$\pAssign{v}$} de conclusie.
\item De laatste afleidingsstap is \RightR{\land}. Dan eindigt $\pi$ in
  \begin{prooftree}
    \AxiomC{}
    \Deduce$\Gamma \fCenter \Delta, !A$
    \AxiomC{}
    \Deduce$\Gamma \fCenter \Delta, !B$
    \RightLabel{\RightR{\land}}
    \BinaryInf$\Gamma \fCenter \Delta, !A \land !B$
  \end{prooftree}
  Beschouw een \iftag{FOL}{!!a{structure}~$\Struct
    M$}{!!a{valuation}~$\pAssign{v}$}. Als
  \iftag{FOL}{$\Struct{M}$}{$\pAssign{v}$} de sequent $\Gamma \Sequent
  \Delta$ vervult, zijn we klaar. Neem dus aan dat dit niet zo is.
  Volgens de inductiehypothese is $\Gamma \fCenter \Delta, !A$ geldig,
  zodat
  $\iftag{FOL}{\Sat{M}}{\pSat{v}}{!A}$. Evenzo is $\Gamma
  \Sequent \Delta, !B$ geldig, zodat
  $\iftag{FOL}{\Sat{M}}{\pSat{v}}{!B}$. Dan
  $\iftag{FOL}{\Sat{M}}{\pSat{v}}{!A \land !B}$.
\item De laatste afleidingsstap is \LeftR{\lor}: opgave.
\item De laatste afleidingsstap is \LeftR{\lif}. Dan eindigt $\pi$ in
  \begin{prooftree}
    \AxiomC{}
    \Deduce$\Gamma \fCenter \Delta, !A$
    \AxiomC{}
    \Deduce$!B, \Pi \fCenter \Lambda$
    \RightLabel{\LeftR{\lif}}
    \BinaryInf$!A \lif !B, \Gamma, \Pi \fCenter \Delta, \Lambda$
  \end{prooftree}
  Beschouw opnieuw
  \iftag{FOL}{!!a{structure}~$\Struct{M}$}{!!a{valuation}~$\pAssign{v}$}
  en neem aan dat \iftag{FOL}{$\Struct{M}$}{$\pAssign{v}$} de sequent
  $\Gamma, \Pi \Sequent \Delta, \Lambda$ niet vervult. We moeten
  aantonen dat $\iftag{FOL}{\Sat/{M}}{\pSat/{v}}{!A \lif !B}$. Omdat
  \iftag{FOL}{$\Struct{M}$}{$\pAssign{v}$} de sequent $\Gamma, \Pi
  \Sequent \Delta, \Lambda$ niet vervult, vervult hij noch $\Gamma
  \Sequent \Delta$, noch $\Pi \Sequent \Lambda$. Omdat $\Gamma
  \Sequent \Delta, !A$ geldig is, geldt
  $\iftag{FOL}{\Sat{M}}{\pSat{v}}{!A}$. Omdat $!B, \Pi \Sequent
  \Lambda$ geldig is, geldt
  $\iftag{FOL}{\Sat/{M}}{\pSat/{v}}{!B}$. Daaruit volgt
  $\iftag{FOL}{\Sat/{M}}{\pSat/{v}}{!A \lif !B}$, zoals verlangd.
\end{enumerate}
\end{proof}

\tagprob{FOL}
\begin{prob}
Voltooi het bewijs van \olref[fol][seq][sou]{thm:sequent-soundness}.
\end{prob}
\tagendprob

\tagprob{notFOL}
\begin{prob}
Voltooi het bewijs van \olref[pl][seq][sou]{thm:sequent-soundness}.
\end{prob}
\tagendprob

\begin{cor}
\ollabel{cor:weak-soundness}
Als $\Proves !A$, dan is $!A$ \iftag{FOL}{geldig}{een tautologie}.
\end{cor}

\begin{cor}
\ollabel{cor:entailment-soundness}
Als $\Gamma \Proves !A$, dan $\Gamma \Entails !A$.
\end{cor}

\begin{proof}
Als $\Gamma \Proves !A$, dan is er voor een eindige deelverzameling
$\Gamma_0 \subseteq \Gamma$ !!a{derivation} van $\Gamma_0 \Sequent
!A$. Volgens \olref{thm:sequent-soundness} maakt elke
\iftag{FOL}{!!{structure}~$\Struct{M}$}{!!{valuation}~$\pAssign{v}$}
ofwel een $!B \in \Gamma_0$ onwaar, ofwel $!A$ waar. Als dus
$\iftag{FOL}{\Sat{M}}{\pSat{v}}{\Gamma}$, dan ook
$\iftag{FOL}{\Sat{M}}{\pSat{v}}{!A}$.
\end{proof}

\begin{cor}
\ollabel{cor:consistency-soundness}
Als $\Gamma$ vervulbaar is, dan is hij consistent.
\end{cor}

\begin{proof}
We bewijzen de contrapositie. Stel dat $\Gamma$ niet consistent is.
Dan bestaan er een eindige $\Gamma_0 \subseteq \Gamma$ en
!!a{derivation} van $\Gamma_0 \Sequent \quad$. Volgens
\olref{thm:sequent-soundness} is $\Gamma_0 \Sequent \quad$ geldig.
Met andere woorden: voor elke
\iftag{FOL}{!!{structure}~$\Struct{M}$}{!!{valuation}~$\pAssign{v}$},
is er een $!C \in \Gamma_0$ waarvoor
$\iftag{FOL}{\Sat/{M}}{\pSat/{v}}{!C}$. Omdat $\Gamma_0 \subseteq
\Gamma$, behoort die $!C$ ook tot~$\Gamma$. Geen enkele
\iftag{FOL}{$\Struct{M}$}{$\pAssign{v}$} vervult dus~$\Gamma$, zodat
$\Gamma$ niet vervulbaar is.
\end{proof}
\end{document}
